\section{Inferencia estadística y estudios complementarios}\label{sec:inferencia}

\subsection{Contrastes pareados, pruebas múltiples y dependencia temporal}

Para comparar dos procedimientos $m_A$ y $m_B$ sobre las series alineadas fuera de muestra (agregadas por fecha en $D+1$ y por origen en los esquemas multihorizonte), definimos la serie temporal de diferenciales de pérdida $d_t := L_{m_A, t} - L_{m_B, t}$ para $t = 1,\dots,T$, donde $L_{m,t}$ es el error absoluto medio en la fecha u origen $t$ (valores negativos favorecen a $m_A$). Ante la autocorrelación serial propia de pronósticos diarios y de horizontes solapados $H$, la varianza de largo plazo se estima mediante el estimador consistente ante heterocedasticidad y autocorrelación (HAC) de Newey--West \cite{newey1987} con núcleo de Bartlett hasta el rezago $L = \max\{1, \lfloor 4(T/100)^{2/9} \rfloor\}$:
\begin{equation}
\hat{\Omega}_{\mathrm{HAC}} := \hat{\gamma}_0 + 2\sum_{k=1}^{L}\left(1 - \frac{k}{L+1}\right)\hat{\gamma}_k, \qquad
\hat{\gamma}_k := \frac{1}{T}\sum_{t=k+1}^{T}(d_t - \bar{d})(d_{t-k} - \bar{d}).
\label{eq:hac}
\end{equation}
A partir de esta estructura de autocorrelación, el tamaño muestral efectivo se cuantifica como
\begin{equation}
T_{\mathrm{eff}} := \frac{T}{\max\!\left\{1,\; 1 + 2\sum_{k=1}^{L}\left(1 - \frac{k}{L+1}\right)\frac{\hat{\gamma}_k}{\hat{\gamma}_0}\right\}}.
\end{equation}
Cuando $T_{\mathrm{eff}} < 2H$ (frecuente en horizontes $H=21$ y $H=28$ días con orígenes no solapados), el escenario se etiqueta automáticamente como \texttt{T\_EFF\_INSUFICIENTE} y se excluye de las afirmaciones inferenciales confirmatorias.

Sobre los escenarios con $T_{\mathrm{eff}} \ge 2H$, el protocolo aplica cinco procedimientos complementarios implementados en \texttt{evaluacion\_estadistica.py} e \texttt{inference.py}:
\begin{enumerate}
    \item \textbf{Prueba DM--HLN--HAC:} estadístico de Diebold--Mariano \cite{diebold1995} con varianza \eqref{eq:hac} y corrección de muestra finita de Harvey--Leybourne--Newbold \cite{harvey1997}, evaluado frente a la distribución $t_{T-1}$ bilateral bajo $H_0 : \mathbb{E}[d_t] = 0$.
    \item \textbf{Model Confidence Set (MCS):} procedimiento de eliminación secuencial de Hansen, Lunde y Nason \cite{hansen2011mcs} al nivel $\alpha = 0{,}10$ con estadístico $T_{\max}$ sobre las pérdidas relativas de los $23$ modelos simultáneos y $B=500$ réplicas de bootstrap estacionario \cite{politis1994stationary} con longitud media de bloque $7$.
    \item \textbf{Prueba global SPA:} contraste de capacidad predictiva superior de Hansen \cite{hansen2005spa} con recienteo consistente y bootstrap estacionario ($B=500$), evaluando si algún competidor supera a una referencia dada sin sesgo de \textit{data snooping}.
    \item \textbf{Prueba de Giacomini--White (GW) auditada:} contraste de capacidad predictiva condicional \cite{giacomini2006tests} con instrumentos $(1, d_{t-1})$ y matriz HAC. Durante la auditoría de calibración bajo $H_0$ ($N_{\mathrm{MC}} = 120$, $T=120$, errores $\mathrm{AR}(1)$ con $\rho=0{,}4$), se detectó que con $\tau=1$ ($L=0$) GW sobrerrechaza ($0{,}2083$ frente a $0{,}05$ nominal) si $|e_t|$ presenta autocorrelación residual, mientras que DM--HLN--HAC ($0{,}0167$) y MCS ($0{,}9333$ de supervivencia) permanecen bien calibrados. Por ello, \texttt{giacomini\_white\_audited} mide la autocorrelación de primer orden $\hat{\rho}_1$ y excluye a GW como prueba confirmatoria cuando $\tau=1$ y $|\hat{\rho}_1| > 0{,}15$.
    \item \textbf{Prueba de Clark--West (CW):} reservada estrictamente para comparaciones entre modelos anidados \cite{clark2007nested} (como SM-CFHR frente a variantes simplificadas o \texttt{MM30}), ajustando la varianza por el sesgo de estimación bajo la nula.
\end{enumerate}

Para controlar la multiplicidad de contrastes, cada familia preespecificada $(\text{suite}, \text{referencia})$ aplica la corrección secuencial de Holm--Bonferroni \cite{holm1979} (control de FWER) y el ajuste de Benjamini--Hochberg \cite{benjamini1995fdr} (control de FDR) únicamente sobre las filas válidas con $T_{\mathrm{eff}} \ge 2H$. La magnitud práctica de las diferencias se reporta mediante la diferencia media $\bar{d}$ con intervalo de confianza al $95\%$ por bootstrap estacionario ($1\,000$ réplicas), la $d$ de Cohen pareada $\bar{d}/s_d$ y el delta de Cliff no paramétrico \cite{cliff1993dominance}. Finalmente, debe recordarse que la agrupación espacial consta de solo $G=7$ conglomerados, por lo que cualquier extrapolación espacial debe interpretarse con las cautelas propias de pocos conglomerados \cite{cameron2015cluster}.

\subsection{Regla formal de decisión para el criterio R3}

Para evitar interpretaciones optimistas de resultados mixtos, el módulo \texttt{evaluate\_r3\_superiority} impone una regla de decisión literal de cuatro condiciones conjuntas para acreditar superioridad empírica (criterio R3):
\begin{enumerate}
    \item \textbf{C1 (Métrica externa):} alcanzar menor $\mathrm{MAE}_{\mathrm{macro}}$ que el mejor modelo de biblioteca en los $28$ escenarios (\texttt{pct\_vs\_MejorLib < 0}).
    \item \textbf{C2 (Tres pruebas distintas calibradas):} acreditar significancia estadística en la mayoría de los escenarios válidos ($T_{\mathrm{eff}} \ge 2H$) mediante tres procedimientos distintos y calibrados: DM--HLN--HAC con corrección de Holm ($p_{\mathrm{Holm}} < 0{,}05$), supervivencia exclusiva en MCS ($\alpha = 0{,}10$) y rechazo en la prueba global SPA ($p_{\mathrm{SPA}} < 0{,}05$), excluyendo GW si hay advertencia de autocorrelación.
    \item \textbf{C3 (Tamaño del efecto):} alcanzar $|d_{\mathrm{Cohen}}| \ge 0{,}20$ (o $|\delta_{\mathrm{Cliff}}| \ge 0{,}147$) en al menos la mitad de los escenarios válidos.
    \item \textbf{C4 (Alcance y consistencia):} consistencia multirrégimen junto con autorización explícita del alcance exclusivamente solar sin conjuntos externos.
\end{enumerate}
Si falta cualquiera de estas cuatro condiciones tras integrar los artefactos definitivos, el veredicto formal del manuscrito debe declararse como \textbf{NO DEMOSTRADO} (o \texttt{NO\_CUMPLE}), detallando con transparencia los regímenes específicos donde SM-CFHR aporta ventaja o resulta superado.

\subsection{Protocolo de estudios complementarios post-OOF}

Una vez persistidas las predicciones definitivas, el flujo \texttt{run\_post\_oof\_stages} ejecuta automáticamente (en modo \texttt{CONFIRMATORY}, prohibiendo todo respaldo sintético) seis estudios complementarios implementados en \texttt{robustness.py}:
\begin{itemize}
    \item \textbf{Ablación controlada ($6$ variantes):} compara \texttt{SM-CFHR\_completo} frente a \texttt{sin\_huber\_L2} ($\delta = 10^6$), \texttt{sin\_l2\_estructurada} ($\alpha_{\mathrm{meteo}} = \alpha_{\mathrm{lin}}$), \texttt{sin\_monotonia} (\texttt{use\_mono=False}), \texttt{escala\_std\_vs\_mad} (estandarización de predictores por media/desviación estándar externa y nota de correspondencia en código) y \texttt{energia\_bruta\_vs\_cf} (ajuste directo en kWh).
    \item \textbf{Sensibilidad 1D y estabilidad por semillas:} barridos unidimensionales de $\delta \in \{0{,}5, 1{,}0, 1{,}35, 2{,}0, 3{,}0\}$, $\alpha_{\mathrm{lin}} \in \{10^{-4},\dots,1\}$ y $\alpha_{\mathrm{meteo}} \in \{10^{-3},\dots,10\}$, junto con evaluación en $10$ semillas ($2026$--$2035$) para verificar determinismo de Fixed y variabilidad de HPO.
    \item \textbf{Estrés con atípicos y ruido:} inyección de $0\%$, $5\%$ y $10\%$ de perturbaciones severas en la respuesta ($\pm 0{,}6$ en factor de capacidad) y ruido gaussiano ($\sigma = 0{,}10$) en predictores, comparando Huber ($\delta=1{,}35$) frente a la referencia $L_2$ ($\delta=10^6$).
    \item \textbf{Curvas de aprendizaje y escalabilidad:} entrenamiento en fracciones $\{10\%, 25\%, 50\%, 100\%\}$ midiendo $\mathrm{MAE}$, tiempo de ajuste, tiempo de predicción y pico de memoria trazada.
    \item \textbf{Incertidumbre conforme temporal:} construcción de intervalos de predicción al $80\%$ y $90\%$ nominales ($\alpha \in \{0{,}20, 0{,}10\}$) mediante predicción conforme dividida temporal sobre residuos absolutos de calibración con corrección finita $q_{1-\alpha} = \mathrm{cuantil}\big(|r_{\mathrm{cal}}|,\; \min\{1, (1-\alpha)(1+1/n_{\mathrm{cal}})\}\big)$ \cite{vovk2005conformal,lei2018conformal}, reportando cobertura global, por planta y ancho medio, bajo la cautela de posible deriva temporal más allá de la intercambiabilidad \cite{barber2023nonexchangeable}.
    \item \textbf{Interpretabilidad aditiva exacta:} dado que SM-CFHR es lineal en el espacio estandarizado $Z$, la contribución exacta de cada predictor respecto a la media de entrenamiento viene dada por $\phi_{q,j} = \sigma_y (z_{q,j} - \bar{z}_{\mathrm{train},j})\,\hat{w}_j$, satisfaciendo $\hat{b}_{\mathrm{fis}} + \sum_{j=1}^p \phi_{q,j}$ igual a la predicción no truncada, sin requerir aproximaciones heurísticas.
\end{itemize}
